% Lab 5 -- Exercise 6  |  Nguyễn Huỳnh Hữu Đức – 2510290  |  Original: Ch.6 3.2 (Enhancement-mode MOSFET)  |  Status: Adjust
% Netlist:
%   VDS d 0 10 ; VGS g 0 0 ; M1 d g 0 0 EM l=1u w=1u     (symbol nmos4)
\exercise{Enhancement-type MOSFET}{Adjust}{Ch.6 3.2 -- Enhancement-mode MOSFET validation}
\statusreason{The original section is a demonstration with the PSpice part \texttt{MbreakN}. It becomes an exercise in
which students extract $V_{GS(\mathrm{Th})}$ and $k$ from the transfer curve; the device is defined with an explicit
level-1 \texttt{.model} so that the result can be compared with theory.}

\task{An n-channel enhancement-type MOSFET is biased with $V_{DS} = 10$\,V. Sweep $V_{GS}$ from $0$ to $6$\,V and plot
$I_D(V_{GS})$. From the curve, determine the threshold voltage $V_{GS(\mathrm{Th})}$ and the constant $k$ of
$I_D = k\,(V_{GS} - V_{GS(\mathrm{Th})})^2$.}

\subsubsection*{Circuit}
\circuitfig{ex06}{Circuit of Exercise 5.6: E-MOSFET transfer characteristic}
\components{
\begin{itemize}[leftmargin=*,itemsep=1pt]
  \item $V_{DS} = 10$\,V drain to ground, $V_{GS}$ gate to ground (swept), source and bulk grounded.
  \item M1 (symbol \texttt{nmos4}, $L = W = 1\,\mu$m): \texttt{.model EM NMOS(Level=1 Vto=2 Kp=1m Lambda=0)}
\end{itemize}}

\subsubsection*{Theoretical Calculation}
\begin{enumerate}[itemsep=2pt]
  \item From the model, write $V_{GS(\mathrm{Th})}$ and $k = \tfrac{K_P}{2}\tfrac{W}{L}$.
  \item Compute $I_D$ at $V_{GS} = 3$, $4$, $5$ and $6$\,V.
  \item Explain how $k$ can be obtained from one measured point $(V_{GS}, I_D)$ once $V_{GS(\mathrm{Th})}$ is known.
\end{enumerate}
\answerspace{8}
\begin{solution}
$V_{GS(\mathrm{Th})} = 2$\,V, $k = 0.5\,\mathrm{mA/V^2}$.
\begin{center}
\begin{tabular}{c c c c c}
\toprule
$V_{GS}$ (V) & 3 & 4 & 5 & 6\\ \midrule
$I_D$ (mA) & 0.5 & 2.0 & 4.5 & 8.0\\
\bottomrule
\end{tabular}
\end{center}
From a measured point: $k = I_D/(V_{GS}-V_{GS(\mathrm{Th})})^2$, e.g.\ $k = 2\,\mathrm{mA}/(2\,\mathrm{V})^2 = 0.5\,\mathrm{mA/V^2}$.
\end{solution}

\subsubsection*{Simulation}
\begin{enumerate}[itemsep=2pt]
  \item Run the DC sweep and plot \texttt{Id(M1)}.
  \item Read $V_{GS}$ where $I_D$ starts to rise and $I_D$ at $V_{GS} = 4$\,V and $6$\,V with the cursor.
\end{enumerate}
\begin{spice}
.model EM NMOS(Level=1 Vto=2 Kp=1m Lambda=0)
.dc VGS 0 6 0.01
\end{spice}
\simfig{ex06_sim}{Transfer characteristic $I_D(V_{GS})$ of Exercise 5.6}

\subsubsection*{Theory vs. Simulation}
\begin{comparison}{Theory vs. simulation, Exercise 5.6}
\cmprow{$V_{GS(\mathrm{Th})}$}{V}{2.00}{}{}
\cmprow{$I_D$ at $V_{GS}=4$\,V}{mA}{2.00}{}{}
\cmprow{$I_D$ at $V_{GS}=6$\,V}{mA}{8.00}{}{}
\cmprow{$k$}{mA/V$^2$}{0.50}{}{}
\end{comparison}
\answerspace{3}
\begin{solution}
Reading the threshold from the knee of the curve is less precise than computing it from two points; both methods can be
discussed. \todo{Fill in the simulation column from the LTspice run.}
\end{solution}
